\anexo{Orientações Gerais sobre Anexos}
\label{an:orientacoes-gerais-sobre-anexos}

Elemento opcional (Guia UECE 2026, 2.2.4.4). Anexo é um texto ou documento que \textbf{não} foi elaborado pelo autor do TCC -- ele serve apenas de fundamentação, comprovação ou ilustração para o trabalho (ex.: uma norma técnica, um relatório de terceiros, uma tabela publicada por um órgão oficial). Se o documento foi produzido por você, ele é Apêndice, não Anexo (ver \autoref{ap:orientacoes-gerais-sobre-apendices}).

As mesmas regras de identificação dos apêndices se aplicam aqui: letras maiúsculas consecutivas seguidas de travessão e título (ANEXO A -- ..., ANEXO B -- ...), com letras dobradas (AA, AB...) quando o alfabeto se esgotar; e a palavra ANEXO, no Sumário e no corpo do texto, deve aparecer em maiúsculas e negrito.

Para acrescentar um novo anexo, crie um arquivo .tex nesta pasta e adicione a linha \texttt{\textbackslash input\{...\}} correspondente na seção \texttt{\textbackslash imprimiranexos} do documento.tex. Se o seu trabalho não precisar de anexos, apague os arquivos desta pasta e remova a linha \texttt{\textbackslash imprimiranexos}.
